\newcommand{\StartDate}{2021-10-11}
\newcommand{\EndDate}{2026-10-08}
\newcommand{\PriceDays}{1254}
\newcommand{\ReturnDays}{1253}
\newcommand{\SplitDate}{2024-10-11}
\newcommand{\InDays}{755}
\newcommand{\OutDays}{498}
\newcommand{\FullDen}{2.591}
\newcommand{\FullInterpretation}{The unrestricted tangency has 94.59\% fitted log-proxy return, 59.46\% volatility and gross exposure 16.93 times equity. Its largest position is 683.31\% in SPY. These weights require substantial shorting and are not an investable recommendation. The long-only tangency instead combines 44.75\% SPY and 55.25\% GLD, with proxy Sharpe 0.868. The long-only MVP has volatility 5.90\%, versus 3.45\% unrestricted.}
\newcommand{\ISDiagnostic}{On the training sample, $d=-2.475<0$. Thus there is no finite fully invested unrestricted Sharpe maximum. Its supremum is $\sqrt{C-B^2/A}=1.901$, approached only with unbounded leverage; the normalized closed form is a Sharpe minimum, not a tangency maximum. The extra bounded diagnostic imposes $|w_i|\le3$ and solves a different problem. It is labeled separately and does not replace the unavailable unrestricted optimum.}
\newcommand{\OOSInterpretation}{The long-only tangency earns Sharpe 1.052 versus 0.796 for equal weight, with total gains of 51.00\% and 25.57\%, respectively. Its maximum drawdown is -17.50\% versus -10.54\%. The bounded diagnostic earns 1.060 Sharpe but draws down -46.32\%. Both MVPs reduce volatility but underperform the benchmark in Sharpe. Higher realized return therefore does not imply safer implementation.}
